\chapter{Open-Source-Veröffentlichung}
\label{cha:opensource}

\section*{Motivation}

Die Entwicklung der HydroNodeStation01 folgte von Beginn an einer Überzeugung: Nützliche Technologie entfaltet ihre volle Wirkung erst dann, wenn sie frei zugänglich ist. Hyperlokal aufgelöste Umweltdaten sind ein gesellschaftliches Gut — je mehr Menschen Zugang zu dieser Technologie haben, desto dichter, aussagekräftiger und demokratischer werden die entstehenden Messnetzwerke.

Fertige kommerzielle Wetterstationen sind oft teuer, geschlossen und nicht auf individuelle Anforderungen anpassbar. Open-Source-Alternativen hingegen ermöglichen es Forschern, Citizen-Scientists und Entwicklern, das Design zu verstehen, zu adaptieren und weiterzuentwickeln. Wer eigene Sensoren hinzufügen, das Sendeintervall anpassen oder das PCB-Layout für seine Anforderungen modifizieren möchte, kann das mit einer offenen Plattform tun — ohne proprietäre Werkzeuge, Lizenzgebühren oder Abhängigkeit von einem Hersteller.

Die modulare I2C-Architektur der HydroNodeStation01 ist bewusst auf dieses Ziel ausgelegt: Jeder kompatible Sensor lässt sich hinzufügen, ohne die Kernarchitektur anzutasten. Die Firmware ist klar strukturiert und soll nach der Veröffentlichung so generalisiert werden, dass sie mit minimalem Aufwand für beliebige Sensorkonstellationen konfiguriert werden kann.

\section*{OSWHLAB Start 2026 Contest}

Das Projekt HydroNodeStation01 wurde beim \textbf{OSWHLAB Start 2026 Open Source Contest} angemeldet — einer Initiative, die Open-Source-Hardware-Projekte fördert, ihnen Sichtbarkeit verleiht und die Autoren mit Sachleistungen für die Fertigung unterstützt. Im Rahmen dieser Förderung wurde ein \textbf{Gutschein im Wert von 150~USD} für die JLCPCB-Fertigung bereitgestellt, der einen substanziellen Teil der Produktionskosten (Gesamtkosten inkl.\ Bestückung, Versand und Steuern: ca.\ 200\,€) abdeckte.

Der finanzielle Aspekt war jedoch nicht die primäre Motivation für die Teilnahme. Der OSWHLAB-Contest legt Wert auf \textit{echte} Open-Source-Konformität: Schaltpläne, PCB-Layouts und Firmware müssen vollständig, reproduzierbar und ohne proprietäre Werkzeuge zugänglich sein. Dieser Anspruch deckt sich exakt mit der Projektphilosophie — die Teilnahme am Wettbewerb war eine logische Konsequenz dieser Ausrichtung, nicht ihr Auslöser.

\section*{Lizenz: CC BY-NC-SA 4.0}

PCB, Firmware (Mikrocontroller-Firmware der HydroNodeStation01) und das Stevenson-Screen-Gehäuse werden unter der \textbf{Creative Commons Attribution-NonCommercial-ShareAlike 4.0 International} (CC BY-NC-SA 4.0) Lizenz veröffentlicht.

Diese Lizenz:
\begin{itemize}
  \item \textbf{erlaubt} das freie Verwenden, Kopieren, Teilen und Anpassen der Materialien für jeden nichtkommerziellen Zweck,
  \item \textbf{verlangt} die Namensnennung des ursprünglichen Autors in allen abgeleiteten Werken,
  \item \textbf{verbietet} die kommerzielle Nutzung ohne separate, schriftliche Genehmigung,
  \item \textbf{verlangt}, dass abgeleitete Werke unter denselben Lizenzbedingungen (CC BY-NC-SA) weitergegeben werden.
\end{itemize}

Diese Wahl stellt sicher, dass die Plattform dauerhaft frei für Forschung, Bildung, Citizen Science und private Projekte bleibt, während eine kommerzielle Vereinnahmung ohne Rücksprache ausgeschlossen ist.

\textit{Hinweis:} Die Lizenz gilt für die stationsspezifische Hardware (PCB, Schaltplan, Gehäuse) und die Mikrocontroller-Firmware (\texttt{stm32\_node}). Das \textit{HydroNode-Backend}, die iOS-App und das HydroNode-Netzwerk unterliegen separaten Lizenzen und sind nicht Gegenstand dieser Veröffentlichung.

\section*{Geplante Veröffentlichung und Generalisierung}

Die Veröffentlichung erfolgt nach erfolgreichem Abschluss des PCB-Tests und der ersten Feldvalidierung. Das Ziel ist eine Plattform, die jemand mit grundlegenden Embedded-Kenntnissen hernehmen, auf seine eigene Sensorausstattung anpassen und als vollständig funktionsfähige, energieeffiziente Messstation in Betrieb nehmen kann.

Geplante Inhalte der Veröffentlichung:

\begin{itemize}
  \item \textbf{Hardware:} Vollständiger Export aller EasyEDA-Pro-Dateien (Schaltplan, PCB-Layout, Gerber-Dateien, BOM) auf GitHub; zusätzlich ein JLCPCB-fertiger Assembly-Export für direkte Nachbestellung. Angestrebte OSHWA-Zertifizierung des Designs.

  \item \textbf{Firmware:} Veröffentlichung auf GitHub inklusive vollständiger Dokumentation des Build-Systems, einer schrittweisen Inbetriebnahme-Anleitung (Commissioning) und einer Beschreibung der Konfigurationsparameter. Die Firmware wird so überarbeitet und generalisiert, dass Sensoren über Konfigurationsdateien aktiviert oder deaktiviert werden können und Sendeintervalle, Batterie-Schwellwerte sowie Payload-Layouts leicht angepasst werden können — ohne tiefen Eingriff in den LoRaWAN-Stack.

  \item \textbf{Gehäuse:} STL-Dateien und Quelldateien (FreeCAD) des angepassten Stevenson-Screen-Gehäuses nach dessen Fertigstellung.

  \item \textbf{Dokumentation:} Diese Studiendokumentation sowie die Energieanalyse (\texttt{power\_budget.md}) als Referenz für eigene Anpassungen.
\end{itemize}

Die Generalisierung der Firmware ist dabei besonders wichtig: Die aktuelle Implementierung ist eng auf die sechs verbauten Sensoren zugeschnitten. Für die Open-Source-Veröffentlichung soll sie so abstrahiert werden, dass beispielsweise ein Nutzer ohne SPS30 oder mit einem anderen CO\textsubscript{2}-Sensor die Plattform ohne Neuschreiben der Kern-Logik einsetzen kann. Die I2C-Architektur und der modulare Sensor-Treiber-Aufbau legen hierfür bereits eine gute Grundlage.
